\chapter*{Summary}\addcontentsline{toc}{chapter}{Summary}

\paragraph{What was reviewed.} The blueprint of the leanVM proof system in leanerVM
(\file{docs/roadmap/protocol-blueprint.md}) and the code built from it on \code{main} at
\code{b435631}: the spine (the abstract instance, the relation, the seams, the composition and
the two master theorems), Layer 0, Layer 1 and the public-input phase. The ground truth is
leanVM at the pin \code{a386121f}: the specification's tex, the Rust verifier (the deployed
protocol) and the Python verifier. Mid-review the owner upgraded \code{main} to Lean 4.34.1 with
new library pins; every statement here is about \code{b435631}, and \cref{sec:map-revision}
says what the upgrade changes (the proof-system statements are the same; the pins moved to
what the review had examined as upstream; a numeral in $\K$ now means something else).

\paragraph{How.} Twelve dossiers, each by an agent with its own context, on the ground truth
phase by phase, on the code, on the libraries, on the boundary with leanISA, on the
documentation and on the literature; every citation of the four ground-truth dossiers checked
against the sources by a second agent (about 1,400 citations; a dozen wrong theorem or line
numbers, no finding weakened); forty-odd Lean and Python probes, re-run at the new pins; a
hostile second reading of eleven major findings that do not concern the opening, which
confirmed nine and weakened two (\cref{sec:faith-rederived}); one register of 138 findings, 115 negative results
and 58 items not verified.

\paragraph{The verdict on the three criteria.}
\begin{description}
\item[Faithfulness.] The blueprint is faithful to leanVM in the bus phase's checks, in the
  table sumcheck's numbers, in the relation's five clauses and in Layer 1's transcriptions,
  which the review confirmed line by line. It is not faithful in the opening (it runs a
  sumcheck leanVM does not run, because it gives the stack an evaluation oracle where leanVM
  commits to an inner-product scheme), in the GKR (a round message of five coefficients where the wire carries four, a layer
  check with an equality factor the deployed one lacks, one combiner too few, and a
  generic component that cannot be appended), in the Flock interface (it drops the eighteen
  limb claims and attributes to the protocol a defect of the Rust prover), in the public-input
  phase as built (it proves the specification's check, which none of the three executable
  verifiers runs), and in the compilation, whose theorems are not well formed and whose error
  lacks three terms.
\item[Non-vacuity.] The relation and the seams are inhabited and refuted as they should be,
  and every extractor computes. But the knowledge-soundness master theorem has content only
  with a bound on the error, and the spine states none: phases that check nothing, at declared
  error one, satisfy its hypotheses for every instance. The check on the public-input message
  is not load-bearing for any theorem, because the verifier pools what it computes rather than
  what it received. And the master theorems constrain the leanISA instance by nothing until
  the adaptor is built, so a weakened instance (two public lines instead of three, a trivial
  Flock predicate, an aliasing layout) is caught by no theorem on \code{main}.
\item[Auditability.] The named knowledge-soundness theorem unfolds to 100 declarations and
  354 lines of statements; twelve proposals, none changing a theorem's meaning, bring it to
  about 80 and 290 (\cref{ch:auditability}).
\end{description}

\paragraph{The findings that matter most} (the register's numbers in parentheses).
\begin{enumerate}
\item The declared error is unconstrained; the soundness theorem says nothing until the spine
  fixes the error in closed form (\code{R5}; machine-checked).
\item The opening phase runs a sumcheck leanVM does not run; an inner-product oracle repairs
  it and makes WHIR the whole opening (\code{R18}; the interface typechecks at both pins).
\item The bus seam admits statements the deployed table sumcheck cannot serve (\code{R2}),
  and the sumcheck's tables are not distinguished from the instance's other tables
  (\code{R11}).
\item The public-input phase proves a check no executable verifier runs (\code{R12}), and the
  check it proves is not load-bearing (\code{R13}); pooling the values sent repairs both, and
  the deployed check's key lemma is proved.
\item The GKR is specified with the wrong wire and the wrong layer check, its per-round error
  is loose, it draws one combiner too few, and its generic component cannot be appended in a
  one-oracle phase (\code{R7}, \code{R9}, \code{R10}; the first weakened by the second reading:
  the blueprint's $5/\abs{\E}$ is a valid, loose bound and its \enquote{check on the cofactor}
  is right).
\item The Flock interface drops the limb claims, is circular with the instance, cannot be
  written over an abstract instance, and files the Rust prover's defect under a soundness error
  (\code{R14} to \code{R17}).
\item The two target theorems are unprovable as sketched: one omits the program's
  well-formedness and attaches a probability to a Boolean; the other is false without a resource
  bound (\code{R27} to \code{R29}). The adaptor's instance is not what the blueprint says it is,
  and the boundary blocks have no bridge lemma (\code{R30}, \code{R31}).
\item The non-interactive error omits grinding, the list-size factor and the hash term; the
  assumed Fiat--Shamir and Merkle interfaces are false as stated or inhabited by nothing, and
  are about other constructions than leanVM's one BLAKE2s map (\code{R20} to \code{R26}).
\item ArkLib's existential notion of round-by-round knowledge soundness carries no knowledge
  content (\code{R1}, critical as a warning about the library, not a defect of a leanerVM
  theorem); the spine's named form is adequate with the extractor read.
\end{enumerate}
The status file's finding that the Python verifier omits the caps is false at the pin and must
not be reported upstream (\code{R38}). Three matters are to be reported to leanVM: the
public-input check differs between the specification and its three verifiers; the Rust Flock
prover fails at one challenge value; several constants and layouts of Flock exist only in the
code.

\paragraph{The owner's question.} If the verifier did not check, through some claim, that the
top lane of both public words is zero, would the theorems catch it? Not the master theorems:
the third public line is data of the instance, and every theorem of the spine and of the phase
holds of an instance with two lines. The adaptor's soundness theorem would, since
\lean{SatisfiedBy} fixes the memory cells to the public words whose top limb is zero by type,
and behind it the constraint-soundness theorem; neither is built. The general rule the review
draws is that a check is load-bearing only if the verifier commits to what it checked, and
that every check of a phase must come with the theorem that refutes its omission
(\cref{ch:drift}).

\paragraph{What the report proposes.} Ten recommendations in priority order
(\cref{sec:opt-recommendations}); the proposed change for every finding, as the passage stands
and as proposed (\cref{ch:changes}); one owner per kind of fact in the documentation, with the
status and the tracker as pointers to the blueprint, names in place of the letter codes and one
index for those that stay (\cref{ch:documentation}, \cref{ch:tracker}, \cref{ch:codeindex});
and the mechanisms that keep the Lean verifier tied to the deployed one as the work proceeds
and the pin moves (\cref{ch:drift}). Nothing has been applied: the blueprint, the status and
the trackers are unchanged, and the drafts wait for the owner's go-ahead.
